\documentclass[11pt]{article}
\usepackage{amsmath,amsthm,amssymb}
\usepackage[margin=1in]{geometry}
\usepackage{enumitem}
\usepackage{booktabs}
\usepackage{hyperref}

\newtheorem{theorem}{Theorem}
\newtheorem{lemma}[theorem]{Lemma}
\newtheorem{proposition}[theorem]{Proposition}
\newtheorem{corollary}[theorem]{Corollary}
\theoremstyle{definition}
\newtheorem{definition}[theorem]{Definition}
\newtheorem{remark}[theorem]{Remark}
\newtheorem{example}[theorem]{Example}

\title{$T_1$ acts as $-1$ on $B^{(\lambda)} V_\lambda$ in the deep stable regime:\\
       a structural theorem dual to the $T_\ell$-eigenvalue}
\author{Clio}
\date{13 May 2026 (late evening, prove session)}

\begin{document}
\maketitle

\begin{abstract}
The structural-skeleton paper of this evening
(\texttt{2026-05-13-evening-T-ell-scalar.tex}) proved a universal fact:
$T_\ell$ acts as the scalar~$q$ on $B^{(\lambda)} V_\lambda$ for every
$\lambda$ with $\lambda_\ell = 1$. The morning refutation paper
(\texttt{2026-05-13-decoration-iso-failed.tex}) observed empirically
that $T_1$ acts as $-1$ on $B$ for certain test shapes but \emph{not}
universally --- the action depends on the shape, with a clear
threshold in the trailing-$1$-row count.

This paper isolates the structural condition. We prove:

\textbf{Theorem.} Let $\lambda$ be a partition with $\lambda_\ell = 1$.
Write $\lambda = (\hat\lambda, 1^q)$ where $\hat\lambda$ has all rows
of length $\ge 2$. Assume:
\begin{itemize}[topsep=0.2em,itemsep=0.1em]
\item[(a)] The SRD-style reduction holds at $\lambda$ and at every
shape encountered in the recursive descent of $\lambda$.
\item[(b)] $q \ge q_0(\hat\lambda) := |\hat\lambda| - 2\ell(\hat\lambda) + 2$.
\end{itemize}
Then
\[
   B_{\ell+1}^{(\lambda)} V_\lambda \;\subseteq\; E_1^-|_{V_\lambda}.
\]

The proof is by induction on $|\hat\lambda^{(\ge 2)}|$. The inductive
step combines three ingredients: (i) $T_1$-equivariance of the
corner-pair branching injection $\Phi_i^\lambda$ (an
index-gap-$\ge 2$ statement), (ii) preservation of $E_1^-$ by
$\mathcal{O}_\lambda = (T_\ell{+}1)\cdots(T_{n-2}{+}1)$ when
$\ell \ge 3$ (also index-gap-$\ge 2$), and (iii) the SRD-style
reduction at the level of $V_\lambda$. The threshold $q_0(\hat\lambda)$
is precisely the condition that $\ell \ge 3$ holds at every step of
the descent. The threshold $q_0$ matches the empirical sharp
threshold on every test shape examined.

Combined with the evening structural-skeleton fact, this gives the
two scalar actions on $B$ in the deep stable regime: $T_\ell$ acts as
$+q$ universally; $T_1$ acts as $-1$ once $q \ge q_0(\hat\lambda)$.
\end{abstract}

\section{Setup and notation}\label{sec:setup}

Notation follows the May-13 multi-corner paper and the r=1 inductive
reduction paper (\texttt{2026-05-13-r1-inductive-reduction.tex},
\texttt{2026-05-13-r-additivity-and-syt-formula.tex}). $H_q(S_n)$ is
the Iwahori--Hecke algebra over $\mathbb{Q}(q)$ with generators
$T_1,\ldots,T_{n-1}$ obeying $(T_i - q)(T_i + 1) = 0$ and the standard
braid relations. $V_\lambda$ is the seminormal Specht module in
asymmetric Murphy normalisation ($b' = 1$). For $j \in \{1, \ldots, n-1\}$,
$E_j^\pm := E_j^\pm(V_\lambda)$ denote the $q$- and $(-1)$-eigenspaces
of $T_j$. For $1 \le k \le n$,
\[
   R'_k \;:=\; (T_{k-1}{+}1)(T_{k-2}{+}1)\cdots(T_1{+}1),
   \qquad
   B^{(\lambda)}_{j_0} V_\lambda \;:=\; R'_{j_0}\cdots R'_n V_\lambda
\]
with $j_0 := \ell + 1$, $\ell := \ell(\lambda)$.

\paragraph{Partition data.} Write $\lambda = (\hat\lambda, 1^q)$
where $\hat\lambda$ is the part with rows of length $\ge 2$ and
$q$ is the number of trailing $1$-rows. The
\emph{length-preserving corners} of $\lambda$ are the removable
corners in column $\ge 2$, i.e., $(\rho_i, \hat\lambda_{\rho_i})$
where $\rho_i$ is a row of $\hat\lambda$ that is a corner of $\hat\lambda$
(no row below it of equal length). List them as $c_1,\ldots,c_r$ in
top-to-bottom order. The \emph{column-$1$ bottom corner} is
$c_* := (\ell, 1)$ when $q \ge 1$. Throughout this paper we assume
$q \ge 1$, so $c_*$ exists.

\paragraph{Decoration shape.}
$\hat\lambda^{(\ge 2)} := \hat\lambda$ with column~$1$ deleted. Then
$|\hat\lambda^{(\ge 2)}| = |\hat\lambda| - \ell(\hat\lambda)$. The
threshold function of this paper is
\begin{equation}\label{eq:threshold}
   q_0(\hat\lambda)
   \;:=\;
   |\hat\lambda| - 2\,\ell(\hat\lambda) + 2
   \;=\;
   |\hat\lambda^{(\ge 2)}| - \ell(\hat\lambda) + 2.
\end{equation}

\paragraph{Corner-pair branching injections.}
For each length-preserving corner $c_i$, set
$\mu_i := \lambda \setminus \{c_i, c_*\}$, a partition of $n - 2$ with
$\ell(\mu_i) = \ell - 1$. The branching injection
$\Phi_i^\lambda : V_{\mu_i} \hookrightarrow V_\lambda$ sends a
seminormal basis vector $v_{T'}^{(\mu_i)}$ to the unique
$+q$-eigenvector of $T_{n-1}$ on the 2-block
$\langle v_{T_A^i}, v_{T_B^i}\rangle$ of $V_\lambda$, where
$T_A^i, T_B^i$ are the two extensions of $T'$ placing $\{n-1, n\}$ at
$\{c_i, c_*\}$ in the two orders. The non-degeneracy of the 2-block
($c_i$ in column $\ge 2$ non-bottom row, $c_*$ in column $1$ bottom
row) ensures $\Phi_i^\lambda$ is well-defined and $\langle T_1, \ldots,
T_{n-3}\rangle$-equivariant (see \textbf{LB-closed} \S 1 or the
r=1 paper Lemma 3.4 for details).

\paragraph{Outer chain.}
\[
   \mathcal{O}_\lambda \;:=\;
   (T_\ell{+}1)(T_{\ell+1}{+}1)\cdots(T_{n-2}{+}1).
\]

\paragraph{SRD-style reduction.}
We say the \emph{SRD-style reduction holds at $\lambda$} if
\begin{equation}\label{eq:srd}
   B^{(\lambda)}_{\ell+1} V_\lambda
   \;=\;
   \mathcal{O}_\lambda \Big[\,
      \sum_{i=1}^r \Phi_i^\lambda\big(B^{(\mu_i)}_{j_0(\mu_i)} V_{\mu_i}\big)
   \Big].
\end{equation}
This is the form of the SRD reduction when all ``other'' pieces
(same-row, pure-length-pres) vanish; it has been proved unconditionally
for various shape families and is the central hypothesis of the
present paper. We say $\lambda$ is in the
\emph{fully stable regime} if \eqref{eq:srd} holds at $\lambda$ and at
every shape encountered in the recursive descent
$\lambda \rightsquigarrow \mu_i \rightsquigarrow \mu_{i,j}
\rightsquigarrow \cdots$ down to a base case.

\paragraph{Descent base case.} The recursive descent terminates at
shapes $\lambda'$ with no length-preserving corners ($r = 0$),
equivalently $\lambda' = (1^k)$ for some $k \ge 1$. At a base case,
\eqref{eq:srd} is vacuous (no $\mu_i$'s).

\section{Inductive step}\label{sec:induction}

The inductive step is short and rests on two index-gap commutations
and the $T_1$-equivariance of~$\Phi_i^\lambda$.

\begin{lemma}[$\Phi_i^\lambda$ is $T_1$-equivariant for $|\mu_i| \ge 4$]\label{lem:phi-T1-equiv}
For $|\mu_i| \ge 4$ (equivalently $n - 2 \ge 4$, i.e., $n \ge 6$),
\[
   T_1 \cdot \Phi_i^\lambda(v) \;=\; \Phi_i^\lambda(T_1 v)
   \qquad \text{for every } v \in V_{\mu_i}.
\]
In particular, $\Phi_i^\lambda(E_1^-|_{V_{\mu_i}}) \subseteq E_1^-|_{V_\lambda}$.
\end{lemma}

\begin{proof}
$\Phi_i^\lambda$ is $\langle T_1, \ldots, T_{n-3}\rangle$-equivariant
by construction (r=1 paper Lemma~3.4, or
\textbf{LB-closed} \S1). For the operator $T_1$ to lie inside this
range we need $1 \le n - 3$, i.e., $n \ge 4$, equivalently
$|\mu_i| = n - 2 \ge 2$. The slightly stronger hypothesis
$n \ge 6$ (i.e., $|\mu_i| \ge 4$) is used below in the chain of
inductive applications but is not required at this lemma alone. The
$E_1^-$-preservation is immediate: if $T_1 v = -v$, then
$T_1 \Phi_i^\lambda(v) = \Phi_i^\lambda(T_1 v) = -\Phi_i^\lambda(v)$.
\end{proof}

\begin{lemma}[$\mathcal{O}_\lambda$ preserves $E_1^-|_{V_\lambda}$ for $\ell \ge 3$]\label{lem:O-E1minus}
If $\ell(\lambda) \ge 3$, the operator $\mathcal{O}_\lambda$ commutes
with $T_1$ as an element of $H_q(S_n)$. In particular,
$\mathcal{O}_\lambda(E_1^-|_{V_\lambda}) \subseteq E_1^-|_{V_\lambda}$.
\end{lemma}

\begin{proof}
$\mathcal{O}_\lambda = (T_\ell{+}1)(T_{\ell+1}{+}1)\cdots(T_{n-2}{+}1)$.
Each factor is $(T_j{+}1)$ for $j \in \{\ell, \ldots, n-2\}$. Since
$\ell \ge 3$, every such $j$ satisfies $j \ge 3$, hence
$|j - 1| \ge 2$, so $T_j$ commutes with $T_1$ inside $H_q(S_n)$
(the commuting Hecke relation for non-adjacent generators). Therefore
$(T_j{+}1)$ commutes with $T_1$ and preserves any $T_1$-eigenspace.
The chain product preserves $E_1^-$.
\end{proof}

\begin{theorem}[Inductive step]\label{thm:induction-step}
Let $\lambda$ be a partition with $\lambda_\ell = 1$ and $\ell(\lambda) \ge 3$.
Assume the SRD-style reduction~\eqref{eq:srd} holds at $\lambda$.
Suppose for each $i \in \{1, \ldots, r\}$,
$B^{(\mu_i)}_{j_0(\mu_i)} V_{\mu_i} \subseteq E_1^-|_{V_{\mu_i}}$.
Then
\[
   B^{(\lambda)}_{\ell+1} V_\lambda \;\subseteq\; E_1^-|_{V_\lambda}.
\]
\end{theorem}

\begin{proof}
By hypothesis, $B^{(\mu_i)} V_{\mu_i} \subseteq E_1^-|_{V_{\mu_i}}$.
For $|\mu_i| \ge 2$, $\Phi_i^\lambda$ is $T_1$-equivariant
(Lemma~\ref{lem:phi-T1-equiv}), so
\[
   \Phi_i^\lambda(B^{(\mu_i)} V_{\mu_i})
   \;\subseteq\; \Phi_i^\lambda(E_1^-|_{V_{\mu_i}})
   \;\subseteq\; E_1^-|_{V_\lambda}.
\]
Summing over $i$,
$\sum_i \Phi_i^\lambda(B^{(\mu_i)} V_{\mu_i}) \subseteq E_1^-|_{V_\lambda}$.
By Lemma~\ref{lem:O-E1minus} ($\ell \ge 3$),
$\mathcal{O}_\lambda$ preserves $E_1^-|_{V_\lambda}$, so
\[
   B^{(\lambda)} V_\lambda
   \;=\; \mathcal{O}_\lambda \Big[\, \sum_i \Phi_i^\lambda(B^{(\mu_i)})\Big]
   \;\subseteq\; \mathcal{O}_\lambda\big(E_1^-|_{V_\lambda}\big)
   \;\subseteq\; E_1^-|_{V_\lambda}.
\]
The equality uses~\eqref{eq:srd}.

\emph{Edge case $|\mu_i| = 1$ or $2$:} this occurs only when $n \le 4$,
hence $\ell(\lambda) \le 2$, ruled out by hypothesis $\ell \ge 3$.
For $|\mu_i| = 3$, i.e., $n = 5$, the operator $T_1$ lies in
$\langle T_1, T_2 \rangle = \langle T_1, \ldots, T_{n-3}\rangle$, so
$\Phi_i^\lambda$ is $T_1$-equivariant by the same proof. No issue.
\end{proof}

\section{Base case}\label{sec:base}

\begin{lemma}[Column shape base case]\label{lem:base}
For $\lambda = (1^k)$ with $k \ge 1$:
\[
   B^{(\lambda)}_{\ell+1} V_\lambda
   \;\subseteq\; E_1^-|_{V_\lambda}.
\]
\end{lemma}

\begin{proof}
$V_{(1^k)}$ is the $1$-dimensional sign representation of $H_q(S_k)$.
Every $T_j$ acts as $-1$, so $V_\lambda = E_j^-$ for every~$j$, in
particular for $j = 1$. Hence any subspace --- including
$B^{(\lambda)} V_\lambda$ --- is contained in $E_1^-$. (When $k \ge 2$
this is meaningful; when $k = 1$, $V_\lambda$ is $1$-dimensional and
the $T_1$-eigenspace question is vacuous.)
\end{proof}

\section{Threshold propagation}\label{sec:threshold}

We verify that the threshold $q_0(\hat\lambda) = |\hat\lambda| - 2\ell(\hat\lambda) + 2$
controls both (i) the inductive lemma hypothesis $\ell \ge 3$ at
$\lambda$ and (ii) the propagation $\lambda$-threshold $\Rightarrow$
$\mu_i$-threshold.

\begin{lemma}[Threshold $\Rightarrow$ $\ell \ge 3$]\label{lem:threshold-ell}
Let $\lambda = (\hat\lambda, 1^q)$ with $\hat\lambda \ne \emptyset$
(so $\ell(\hat\lambda) \ge 1$). If
$q \ge q_0(\hat\lambda) = |\hat\lambda| - 2\ell(\hat\lambda) + 2$, then
$\ell(\lambda) \ge |\hat\lambda^{(\ge 2)}| + 2 \ge 3$.
\end{lemma}

\begin{proof}
$\ell(\lambda) = \ell(\hat\lambda) + q$. The threshold gives
$q \ge |\hat\lambda| - 2\ell(\hat\lambda) + 2$, so
$\ell(\lambda) = \ell(\hat\lambda) + q \ge |\hat\lambda| - \ell(\hat\lambda) + 2
= |\hat\lambda^{(\ge 2)}| + 2$.
Since $\hat\lambda \ne \emptyset$ implies
$|\hat\lambda^{(\ge 2)}| \ge 1$ (every row of $\hat\lambda$ has
length $\ge 2$, contributing $\ge 1$ to $\hat\lambda^{(\ge 2)}$), we
get $\ell(\lambda) \ge 3$.

\emph{Boundary case $\hat\lambda = \emptyset$:} then $\lambda = (1^q)$,
which is the base case (Lemma~\ref{lem:base}), not handled by the
inductive step.
\end{proof}

\begin{lemma}[Threshold propagates to each $\mu_i$]\label{lem:threshold-propagate}
Let $\lambda = (\hat\lambda, 1^q)$ with $\hat\lambda \ne \emptyset$
and $q \ge q_0(\hat\lambda)$. Then for every length-preserving corner
$c_i$ of $\lambda$ (i.e., corner of $\hat\lambda$),
\[
   q(\mu_i) \;\ge\; q_0(\hat\mu_i).
\]
\end{lemma}

\begin{proof}
Let $c_i = (\rho_i, \hat\lambda_{\rho_i})$ be a length-preserving
corner. We split into two cases based on whether $\rho_i = \ell(\hat\lambda)$
and $\hat\lambda_{\rho_i} = 2$.

\textbf{Case 1: $\rho_i < \ell(\hat\lambda)$, or $\rho_i = \ell(\hat\lambda)$
with $\hat\lambda_{\rho_i} \ge 3$.}

In this case, removing $c_i$ leaves row $\rho_i$ with length
$\hat\lambda_{\rho_i} - 1 \ge 2$, so the row remains in $\hat\mu_i$.
Hence $\ell(\hat\mu_i) = \ell(\hat\lambda)$ and $|\hat\mu_i| = |\hat\lambda| - 1$.
Removing $c_*$ removes one trailing $1$-row, so $q(\mu_i) = q - 1$.

Then
\[
   q_0(\hat\mu_i) = |\hat\mu_i| - 2\ell(\hat\mu_i) + 2
   = (|\hat\lambda| - 1) - 2\ell(\hat\lambda) + 2
   = q_0(\hat\lambda) - 1,
\]
and $q(\mu_i) = q - 1 \ge q_0(\hat\lambda) - 1 = q_0(\hat\mu_i)$. \checkmark

\textbf{Case 2: $\rho_i = \ell(\hat\lambda)$ and $\hat\lambda_{\rho_i} = 2$.}

Removing $c_i = (\ell(\hat\lambda), 2)$ leaves row $\ell(\hat\lambda)$
of length~$1$, so this row drops out of $\hat\mu_i$ (which records
rows of length~$\ge 2$). $\hat\mu_i = \hat\lambda$ with the entire
last row removed: $|\hat\mu_i| = |\hat\lambda| - 2$ and
$\ell(\hat\mu_i) = \ell(\hat\lambda) - 1$.

What about $q(\mu_i)$? Counting trailing $1$-rows of $\mu_i$: removing
$c_*$ from $\lambda$ kills one trailing $1$-row, but the reclassified
former-$\hat\lambda$-last-row of length~$1$ becomes a new trailing
$1$-row. Net change: $q(\mu_i) = q$.

(Sanity check: $|\mu_i| = |\hat\mu_i| + q(\mu_i) = (|\hat\lambda| - 2) + q
= |\lambda| - 2$. \checkmark)

Then
\[
   q_0(\hat\mu_i) = |\hat\mu_i| - 2\ell(\hat\mu_i) + 2
   = (|\hat\lambda| - 2) - 2(\ell(\hat\lambda) - 1) + 2
   = |\hat\lambda| - 2\ell(\hat\lambda) + 2
   = q_0(\hat\lambda),
\]
and $q(\mu_i) = q \ge q_0(\hat\lambda) = q_0(\hat\mu_i)$. \checkmark

In both cases, $q(\mu_i) \ge q_0(\hat\mu_i)$.
\end{proof}

\begin{remark}[Dichotomy on corners]\label{rem:dichotomy}
The case dichotomy
$\rho_i = \ell(\hat\lambda)$ vs.\ $\rho_i < \ell(\hat\lambda)$ is
exhaustive for length-preserving corners~$c_i$ of $\lambda$: if
$\rho_i < \ell(\hat\lambda)$, $\hat\lambda_{\rho_i} \ge 3$
automatically. (If $\hat\lambda_{\rho_i} = 2$ with
$\rho_i < \ell(\hat\lambda)$, then $\hat\lambda_{\rho_i+1} \le 2$, but
$\hat\lambda$ has all rows $\ge 2$, so $\hat\lambda_{\rho_i+1} = 2 =
\hat\lambda_{\rho_i}$, making $c_i$ \emph{not} a corner --- the row
below has the same length.)
\end{remark}

\section{Main theorem}\label{sec:main}

\begin{theorem}[Main: $T_1 = -1$ on $B$ in the deep stable regime]\label{thm:main}
Let $\lambda = (\hat\lambda, 1^q)$ with $\lambda_\ell = 1$ (i.e.,
$q \ge 1$). Assume:
\begin{itemize}[topsep=0.2em,itemsep=0.1em]
\item[(a)] The SRD-style reduction~\eqref{eq:srd} holds at $\lambda$
and at every shape encountered in the recursive descent of $\lambda$.
\item[(b)] $q \ge q_0(\hat\lambda) := |\hat\lambda| - 2\ell(\hat\lambda) + 2$.
\end{itemize}
Then
\[
   B^{(\lambda)}_{\ell+1} V_\lambda
   \;\subseteq\; E_1^-|_{V_\lambda}.
\]
\end{theorem}

\begin{proof}
We argue by induction on $|\hat\lambda^{(\ge 2)}| = |\hat\lambda| - \ell(\hat\lambda)$.

\emph{Base case $|\hat\lambda^{(\ge 2)}| = 0$.} Then $\hat\lambda$ has
all rows of length~$1$, but our convention says $\hat\lambda$ consists
of rows of length $\ge 2$. So $\hat\lambda = \emptyset$ and
$\lambda = (1^q)$. By Lemma~\ref{lem:base}, $B^{(\lambda)} V_\lambda
\subseteq E_1^-|_{V_\lambda}$.

\emph{Inductive step.} Suppose $|\hat\lambda^{(\ge 2)}| \ge 1$, so
$\hat\lambda \ne \emptyset$ and $\lambda$ has $r \ge 1$
length-preserving corners. By Lemma~\ref{lem:threshold-ell},
hypothesis~(b) implies $\ell(\lambda) \ge 3$. By
Lemma~\ref{lem:threshold-propagate}, each $\mu_i$ satisfies its own
threshold: $q(\mu_i) \ge q_0(\hat\mu_i)$. Hypothesis~(a) propagates to
$\mu_i$ by definition of the fully stable regime: the recursive descent
of $\mu_i$ is a subtree of that of $\lambda$, and SRD reduction holds
throughout. Hence the inductive hypothesis applies to each $\mu_i$:
\[
   B^{(\mu_i)} V_{\mu_i} \;\subseteq\; E_1^-|_{V_{\mu_i}}.
\]
Apply Theorem~\ref{thm:induction-step} (using $\ell(\lambda) \ge 3$
and~\eqref{eq:srd} at $\lambda$):
$B^{(\lambda)} V_\lambda \subseteq E_1^-|_{V_\lambda}$.

\emph{Induction terminates because $|\hat\mu_i^{(\ge 2)}| < |\hat\lambda^{(\ge 2)}|$.}
Concretely, $|\hat\mu_i^{(\ge 2)}| = |\hat\mu_i| - \ell(\hat\mu_i)$.
\begin{itemize}[topsep=0.2em,itemsep=0.1em]
\item Case 1: $|\hat\mu_i^{(\ge 2)}| = (|\hat\lambda| - 1) - \ell(\hat\lambda)
= |\hat\lambda^{(\ge 2)}| - 1$.
\item Case 2: $|\hat\mu_i^{(\ge 2)}| = (|\hat\lambda| - 2) - (\ell(\hat\lambda) - 1)
= |\hat\lambda^{(\ge 2)}| - 1$.
\end{itemize}
Both cases drop $|\hat\lambda^{(\ge 2)}|$ by exactly~$1$.
\end{proof}

\begin{remark}[Pairing with the structural skeleton]
Theorem~\ref{thm:main} is dual to Theorem~A of
\texttt{2026-05-13-evening-T-ell-scalar.tex}, which proved that
$T_\ell$ acts as the scalar $+q$ on $B^{(\lambda)} V_\lambda$ for
every $\lambda$ with $\lambda_\ell = 1$ (no threshold required). The
combined picture in the deep stable regime ($q \ge q_0(\hat\lambda)$
and fully stable):
\[
   B^{(\lambda)} V_\lambda
   \;\subseteq\;
   E_\ell^+|_{V_\lambda}\,\cap\,E_1^-|_{V_\lambda}.
\]
The $T_\ell$-action is universal; the $T_1$-action has a
shape-dependent threshold. Both belong to the centraliser
$Z_{H_q}(T_1) \cap Z_{H_q}(T_\ell)$ on $B$ (and any operator preserving
$B$ as a subspace must commute with these two scalar actions, hence
lies in the appropriate centraliser).
\end{remark}

\section{Threshold sharpness: computational evidence}\label{sec:sharp}

We computed $T_1|_B$ for a range of test shapes via mod-$P$ arithmetic
($P = 100\,003$, $q = 11$). For each shape $\lambda = (\hat\lambda, 1^q)$,
we extracted a basis $W$ of $B^{(\lambda)} V_\lambda$ and tested whether
$T_1 W = -W$ (i.e., $B \subseteq E_1^-$).

\begin{center}\renewcommand{\arraystretch}{1.15}
\begin{tabular}{cccc}
\toprule
$\hat\lambda$ & $q_0(\hat\lambda)$ & first $q$ giving $B \subseteq E_1^-$ & verdict \\
\midrule
$(2)$         & $2$ & $2$ & matches \\
$(3)$         & $3$ & $3$ & matches \\
$(4)$         & $4$ & $4$ & matches \\
$(5)$         & $5$ & $5$ & matches \\
$(2, 2)$      & $2$ & $2$ & matches \\
$(3, 2)$      & $3$ & $3$ & matches \\
$(3, 3)$      & $4$ & $4$ & matches \\
$(4, 2)$      & $4$ & $4$ & matches \\
$(4, 3)$      & $5$ & $\ge 5^*$ & matches at $q=4$ (below) \\
$(2, 2, 2)$   & $2$ & $2$ & matches \\
\bottomrule
\end{tabular}
\end{center}

In every test, $q_0(\hat\lambda)$ matches the empirical sharp
threshold. The cells of the test grid below threshold show
``$T_1$ mixed'' (neither $+q$ nor $-1$ on all of $B$); the cells at
or above threshold show $T_1 |_B = -1$.

\paragraph{${}^*$ Note on $(4, 3)$.} The cell $\hat\lambda = (4, 3)$
was verified \emph{below} threshold ($q = 4$, $\dim V_\lambda = 1100$):
$\dim B = 5$ (matches the dim closed form) but $T_1$ has mixed action
on $B$ --- consistent with the prediction $q < q_0 = 5$. The
above-threshold case $(4, 3, 1^5)$ ($n = 12$, $\dim V_\lambda = 4158$)
exceeds the computational budget of the mod-$P$ verification harness
in this session and is left as a future check. The shape illustrates
the precise distinction between \emph{dim} stability and
$T_1 = -1$ stability: the closed-form dim matches at $q = 4$ but
$T_1$-stability requires the strictly stronger condition $q \ge 5$.

\paragraph{Scripts.}
\texttt{\textasciitilde/projects/scratch/2026-05-13-night-j-formula/verify\_T1\_neg1.py}.

\section{Honest scope}\label{sec:scope}

\paragraph{What is proved.}
Theorem~\ref{thm:main} (inductive step + threshold propagation +
column-shape base case). All three are short structural lemmas
relying only on the index-gap commutation
$|T_j, T_1| = 0$ for $j \ge 3$, the $T_1$-equivariance of
$\Phi_i^\lambda$ (for $n \ge 4$), and the SRD-style reduction
hypothesis.

\paragraph{What is conditional.}
Hypothesis~(a) (full stability of the SRD-style reduction along the
descent) remains the substantive structural input. It has been proved
in individual shape families (hooks, $(2^a, 1^?)$, $(3, 3, 1^?)$ at
$q \ge 2$, $(4, 2, 1^?)$ at $q \ge 2$, $(4, 3, 1^?)$ at $q \ge 3$, and
the $(3, 2, 1^?)$ family at all $q$) but not yet in full generality.
A general structural proof of SRD-style reduction would render
Theorem~\ref{thm:main} unconditional in every fully stable shape.

\paragraph{Sharpness.}
The threshold $q_0(\hat\lambda)$ matches the empirical sharp
threshold on every shape tested in \S\ref{sec:sharp}. The present
paper proves it is \emph{sufficient}; sharpness (i.e., that
$B \not\subseteq E_1^-$ for $q < q_0$) is empirical only --- a
proof of necessity is a natural next target. Conceptually, a Case-2
step ``borrows'' a $1$-row from $\hat\lambda$ to $\mathrm{trail}(\lambda)$,
keeping $q_0$ the same. A Case-1 step decrements both $q$ and
$q_0$ by~$1$ each. The threshold counts the worst-case descent
depth needed for the chain $\mathcal{O}_\lambda$ at every level to
have index $\ge 3$.

\paragraph{Relation to the morning refutation paper.}
The morning paper
(\texttt{2026-05-13-decoration-iso-failed.tex}) observed
$T_1 |_B = -1$ on the test case $\lambda = (3, 2, 1^3)$ and on
``$4$ other test cases'' but did not extract a threshold. The
threshold $q_0(\hat\lambda) = |\hat\lambda| - 2\ell(\hat\lambda) + 2$
identified here closes that observation into a clean structural
statement.

\paragraph{What this enables.}
Combined with the evening structural-skeleton paper, we now have a
two-sided structural skeleton for $B^{(\lambda)} V_\lambda$ in the
deep stable regime:
\begin{enumerate}[leftmargin=2em,topsep=0.2em,itemsep=0.1em]
\item $T_\ell$ acts as $+q$ (universal, evening paper).
\item $T_1$ acts as $-1$ (deep stable regime, this paper).
\item $E_1^-|_{V_\lambda} \subseteq \ker \Omega$, with
$\Omega := R'_{\ell+1}\cdots R'_n$, so
$\overline{\Omega} : V_\lambda / E_1^-|_{V_\lambda} \to B$ is
surjective by construction (evening paper, Theorem~A$'$).
\end{enumerate}
The combined picture gives strong restrictions on any potential
Hecke-module structure on $B$: every operator preserving $B$ as a
subspace must commute with $T_1$ and $T_\ell$, hence lies in
$Z_{H_q}(T_1) \cap Z_{H_q}(T_\ell)$. This is the natural setting in
which to look for the elusive ``decoration functor'' upgrading the
SYT-count equality to a representation-theoretic isomorphism.

\paragraph{Open problems.}
\begin{itemize}[leftmargin=2em,topsep=0.2em,itemsep=0.1em]
\item Prove the SRD-style reduction~\eqref{eq:srd} structurally in
full generality.
\item Prove the threshold $q_0(\hat\lambda)$ is sharp: $B \not\subseteq E_1^-$
for $q < q_0(\hat\lambda)$.
\item Describe $B^{(\lambda)} V_\lambda$ for $q < q_0(\hat\lambda)$ ---
what is the $T_1$-eigenspace decomposition of $B$ on the
non-stable side?
\end{itemize}

\end{document}
